\section{Research Question 1: Parent-Anchored Behavioral Normalization}
\label{sec:method_rq1}

\subsection{Theoretical Background and Mathematical Formulation}
\label{subsec:rq1_theory}

Due to the lack of a unified coordinate system between architectures, cross architectural backdoor auditing faces difficulties. These can be split into:
\begin{enumerate}
    \item Hidden states have different dimensionality ($d = 768, 1536, 4096$) and Normalisation types (LayerNorm vs. RMSNorm) and Attention types (Multi-Head vs Grouped Query).
    \item Vocabulary size differs ($|V| \in [32000, 128256]$) which causes an arbitrary scaling offset in candidate output logits.
\end{enumerate}

CALB-Shield does not attempt to align the candidate's Weight Tensors. Instead it treats the candidate as a black-box generator and performs a diagnostic analysis using a set of 30 behaviour probes ($\mathcal{P} = \{p_1, \dots, p_{30}\}$) to assess the following five key aspects of the model's semantics: Factual Knowledge, Ethical Reasoning, Technical Analysis, Creative Synthesis and Safety Boundaries.

\subsubsection{180 Dimensional Behavioural Signature}
For each behaviour probe $p_j$, the candidate model generates an unnormalised logit vector $\mathbf{z} \in \mathbb{R}^{|V|}$. This can be converted to a probability distribution $\mathbf{p} = \text{softmax}(\mathbf{z})$ by applying the softmax function. The probability distribution is then ordered in descending order ($p_{(1)} \ge p_{(2)} \ge \dots \ge p_{(|V|)}$). From the distribution, CALB-Shield generates 6 metrics describing the model's behaviour:
\begin{enumerate}
    \item \textbf{Output Token Entropy ($H$):}
    \begin{equation}
    H(p_j) = -\sum_{i=1}^{K} p_{(i)} \ln p_{(i)}
    \end{equation}
    \item \textbf{Top-1 / Top-2 Logit Difference ($\Delta_{\text{logit}}$):}
    \begin{equation}
    \Delta_{\text{logit}}(p_j) = z_{(1)} - z_{(2)}
    \end{equation}
    \item \textbf{Top-5 Probability Mass ($M_5$):}
    \begin{equation}
    M_5(p_j) = \sum_{i=1}^{5} p_{(i)}
    \end{equation}
    \item \textbf{Top-1 Prediction Confidence ($p_1$):}
    \begin{equation}
    p_1(p_j) = p_{(1)}
    \end{equation}
    \item \textbf{Spread ($S$):}
    \begin{equation}
    S(p_j) = \frac{\exp(H(p_j))}{K}
    \end{equation}
    \item \textbf{Mean Top-$K$ Log Probability ($\bar{\ell}$):}
    \begin{equation}
    \bar{\ell}(p_j) = \frac{1}{K} \sum_{i=1}^{K} \ln p_{(i)}
    \end{equation}
\end{enumerate}

The behaviour signature $\mathbf{f}$ is generated by concatenating the metrics together:
\begin{equation}
\mathbf{f} = \big[ \mathbf{f}(p_1)^\top, \mathbf{f}(p_2)^\top, \dots, \mathbf{f}(p_{30})^\top \big]^\top \in \mathbb{R}^{180}
\end{equation}
Where $\mathbf{f}(p_j) = [f_1(p_j), \dots, f_6(p_j)]^\top$.

\subsubsection{Hardware Repeatability and Stability}
To ensure the logit signature is deterministic, the feature extraction process was performed 150 times on Apple Silicon Metal Unified Memory Architecture (MPS) using greedy decoding (\texttt{temperature = 0.0}). The mean Coefficient of Variation (CV) was found to be 0.0446\% and 174/180 features had zero variance at the bit level (0.0000\%). Therefore, the logit signature is stable.

\begin{figure*}[t]
\centering
\safeincludegraphicspdf[width=0.92\textwidth]{fig2_probe_repeatability.pdf}
\caption{Behaviour Probes With Repeated Hardware Executions. 174 out of 180 probed dimensions (96.7\%) had zero variance ($\text{CV} = 0.0000\%$) over 150 physical executions on Apple Silicon Metal with greedy decoding (\texttt{temperature = 0.0}). The mean CV across all dimensions was 0.0446\%. The entropy of the first token output peaked at 5.61\% (above the 5.0\% limit). This was due to unwarmed GPU thread group. The entropy stabilised to a CV less than 0.20\% over all tokens.}
\label{fig:probe_repeatability}
\end{figure*}

\subsection{Parent Anchored Relative Normalization}
\label{subsec:rq1_norm}
Published checkpoints in open-source model repositories specify what their parent architecture was (e.g. \nolinkurl{Qwen2.5-Coder-1.5B-Instruct}). Let $\mathbf{f}_{\text{cand}} \in \mathbb{R}^{180}$ be the behaviour signature of the candidate model. Let $\boldsymbol{\mu}_{\text{parent}} \in \mathbb{R}^{180}$ and $\boldsymbol{\sigma}_{\text{parent}} \in \mathbb{R}^{180}$ be the mean and standard deviation of the clean parent baseline. The Parent Anchored Relative Transformation is:
\begin{equation}
\mathbf{z}_i = \frac{\mathbf{f}_{\text{cand}, i} - \boldsymbol{\mu}_{\text{parent}, i}}{\boldsymbol{\sigma}_{\text{parent}, i} + \varepsilon}, \quad \forall i \in \{1, \dots, 180\}
\end{equation}
The relative transform ensures that any deviation caused by the candidate is compared to the verified parent baseline, thus cancelling the parent architecture baseline offset.

\subsection{Machine Learning and Deep Learning Classifiers Tested}
\label{subsec:rq1_classifiers}
The following 5 types of classifiers were trained using clean parent anchored baselines from Meta-LLaMA-3-8B and Mistral-7B, to perform zero-shot cross architectural transfer on physical Qwen-2.5-Coder-1.5B checkpoints:
\begin{enumerate}
    \item \textbf{Logistic Regression (L2 Regularized, $C = 1.0$):} A linear classifier with a sigmoid output probability.
    \item \textbf{Linear SVM ($C = 1.0$):} A maximum margin classifier.
    \item \textbf{Random Forest (100 trees, Max Depth = 8):} An ensemble classifier partitioning orthogonal feature spaces.
    \item \textbf{Multi-Layer Perceptron (Neural Network):} Fully connected network with ReLU activation functions and Adam optimiser with dropout ($p = 0.1$) ($180 \to 128 \to 64 \to 1$).
    \item \textbf{Extra Trees Classifier (100 estimators):} Decision tree classifier considering random orthogonal splitting thresholds.
\end{enumerate}

\noindent\textbf{Performance Differences Between Raw and Anchored Features:}
\begin{itemize}
    \item \textbf{A. Raw/Uncalibrated Features:} All 5 classifiers failed during cross architectural transfer, giving only 50.0\% accuracy ($\text{TP} = 0, \text{FP} = 0, \text{TN} = 1, \text{FN} = 1$, 100.0\% False Negative Rate). The Clean Mistral training checkpoint had high natural logit difference ($\Delta_{\text{logit}} = 3.65$). Consequently, the compressed logit difference of the backdoored Qwen model ($\Delta_{\text{logit}} = 3.207$) fell within the normal range and could not be detected.
    \item \textbf{B. Parent Anchored Normalization:} The clean parent anchor was subtracted from the candidate model's difference. This gave a clear relative anomaly ($z \ge +4.12\sigma$) due to the small absolute difference. All classifiers achieved 100.0\% Detection Accuracy on the physical testbed.
\end{itemize}

\subsection{Empirical Evaluation of 9 Zero Reference Approaches}
\label{subsec:rq1_negative}
An ideal goal in model governance is to perform zero reference audit. This means auditing an arbitrary checkpoint without knowing its parent architecture. To determine if this is possible, 9 distinct zero reference detection approaches were implemented across all 4 checkpoints (Clean LLaMA-3-8B, Clean Mistral-7B, Clean Qwen-1.5B, Poisoned Qwen-1.5B PoC).

All 9 uncalibrated approaches failed to detect the candidate:
\begin{enumerate}
    \item \textbf{Static Inter-Model Scaling:} Variation between logit scaling caused trojan differences to be masked. Clean Qwen was identified as an anomaly ($S = 2.028$) due to different baseline scaling.
    \item \textbf{Dynamic Semantic Probing (STRIP Style):} Trojans are inactive on neutral text (Clean Qwen JS = 0.43, Trojan Qwen JS = 0.39). No difference detected.
    \item \textbf{Uncalibrated Shortcut Inversion (BAIT Style):} Pretrained instruction bias caused more sharp universal shortcuts within Clean Mistral (UAS = 2.67) compared to backdoored Qwen (UAS = 1.45). This caused the detector to be inverted.
    \item \textbf{Residual Latent Geometry SVD (BackdoorID Style):} Differences in dimensionality and depth of layers between different architectures made it impossible to find a universal threshold of singular values. Clean Qwen had the same residual shift as the trojan model ($\Delta \sigma_1 = 0.042$).
    \item \textbf{Zero Shot Logit Kurtosis:} Logit tail was related to vocabulary size ($|V| = 152k$ vs $32k$). No relation to trojan presence.
    \item \textbf{Cross Architectural Negative Constraints:} Autoregressive attention caused models to focus on forbidden concepts when told not to (Pink elephant effect). This caused constraint probing to fail.
    \item \textbf{Output Sequence Locking (ConfGuard Style):} Inactive backdoors did not cause high confidence lock out on neutral inputs (Max consecutive = 4). Clean models often locked (13 tokens).
    \item \textbf{Memory Extraction Scanning:} Clean formatting caused false confidence spikes. Subtle trojans did not leak as they required exact trigger syntax.
    \item \textbf{Direct Weight Spectral Scanning (PEFTGuard Style):} SVD performed on full base weight matrices caused the perturbations to be spread out across 48 layers ($\Delta W$ Top-1 energy = 0.0513). This meant that low rank detection failed (AUC = 0.509).
\end{enumerate}

This corroborates the theoretical proof given in~\cite{cispa2026datafree} by CISPA Helmholtz Center for Information Security: uncalibrated data free backdoor detection on modern language models is ill-posed.

\subsection{Upstream Anchored Admission Control}
\label{subsec:rq1_upstream}
Due to the failure in zero reference approaches, it is important to note that models will never be deployed without any knowledge of their parent. Therefore, the parent information can be used in CALB-Shield via Upstream Anchored Admission Control:
\begin{enumerate}
    \item When an untrusted fine tuned model is submitted for admission, the 30 behaviour probes are applied to both the candidate and the declared parent base model.
    \item The admission engine calculates the behavioural difference between models:
    \begin{equation}
    \Delta \mathbf{f} = \mathbf{f}(M_{\text{candidate}}) - \mathbf{f}(M_{\text{parent}})
    \end{equation}
    \item Relative normalisation is applied to cancel out any variation due to the architecture. This brings all models into a common space.
\end{enumerate}

To ensure that parent anchoring can distinguish between legitimate fine tuning changes and malicious backdoors, 60 synthetic feature distributions were tested (Mild task tuning, domain adaptation, DPO alignment and parameter drifting with $\sigma \le 0.35$).

Along with the trojan collapses using 3-Way Leave-One-Architecture-Out (LOAO) cross validation:
\begin{itemize}
    \item \textbf{Hold Out Qwen-1.5B (Linear SVM):} 100.0\% Overall Accuracy, 0.0\% Benign FAR, 0.0\% FNR. The Trojan PoC was quarantined with 100\% accuracy.
    \item \textbf{Hold Out Mistral-7B (Linear SVM):} 100.0\% Overall Accuracy, 0.0\% Benign FAR, 0.0\% FNR.
    \item \textbf{Hold Out LLaMA-3-8B (Linear SVM):} 100.0\% Overall Accuracy, 0.0\% Benign FAR, 0.0\% FNR.
\end{itemize}
Legitimate task adaptation caused the behaviour to be distributed and resulted in high entropy ($H > 0.35$) meaning the models are linearly separable from the trojan collapse.